\section{总结、局限与后续工作}

\subsection{对研究问题的回答}

\textbf{第一，位置变化会不会明显影响分类？}会，而且影响远大于中心测试给人的印象。中心训练的 MLP、CNN 和 ViT-AbsPos 在 Center 上都有 $86\%$ 以上准确率，CNN 达到 $90.93\%$；换成 Large Shift 后，三者只剩 $14.22\%$、$13.80\%$ 和 $20.03\%$。这说明原任务中的高准确率很大一部分建立在“衣物总在画布中央”这一稳定条件上。一旦位置变化并伴随边缘裁剪，模型原有的分类能力不能直接迁移过去。

\textbf{第二，使用绝对位置编码的 Tiny ViT 能否自然适应位移？}不能。ViT-AbsPos 在 Large Shift 和位置网格上比中心训练的 MLP、CNN 略好，但 Grid Accuracy 仍只有约 $21\%$，Robust Drop 达到 66.35 个百分点。受控三种子实验中，固定二维 Sin-Cos 与可学习绝对位置编码的位移表现相近；完全移除位置编码虽然把 Large Shift Accuracy 从 $20.86\%$ 提高到 $39.42\%$，也同时降低中心准确率，而且仍明显不如平移增强。自注意力和位置编码选择都会影响取舍，却不能替代合适的训练分布。

\textbf{第三，数据增强、Mean Pooling 和卷积 stem 是否有效？}三者的作用大小不同。加入平移、旋转和随机擦除后，ViT 的 Grid Accuracy 从约 $21\%$ 提升到约 $65\%$，这是最主要的变化；三种子逐项消融进一步表明，平移负责消除主要位置落差，旋转负责补上角度鲁棒性，而 Random Erasing 在当前设置下基本中性。Mean Pooling 又带来约 1 个百分点的小幅改善；加入卷积 stem 后，HybridConv-ViT 的 Center、Large Shift、Rotation、Shift+Rotation 和 Grid Accuracy 都达到约 $81\%$。它没有取得最高的中心准确率，但位置改变、角度改变或二者同时出现时，成绩几乎不再下降。

因此，可以认为 HybridConv-ViT 已经达到本项目对“位置可变 FashionMNIST 分类器”的初步设计要求：它在保留较好分类准确率的同时，明显减弱了模型对中心位置和固定朝向的依赖，并且这一结果在五个随机种子下保持一致。它还不是一个完全解决问题的模型，例如 Shirt 等相似上装仍较难分类，极端位置的前景裁剪也没有消失；但相较于三个中心训练基线，它已经从“中心图像分类器”推进成了能够实际处理位置与角度变化的分类器。

\subsection{现有局限}

\begin{enumerate}
  \item \textbf{补充消融使用三个随机种子。}四阶段增强消融和位置编码消融均报告 seed 42、2026、3407 的均值与样本标准差，已能检查结论是否依赖单次运行，但重复次数仍少于六组主实验的五个种子；很小的均值差异不宜解释为稳定排序。
  \item \textbf{边缘裁剪仍构成混杂因素。}第 7 节除原有内外圈汇总外，还在完全相同的欧氏位移半径上比较无裁剪与裁剪坐标，排除了“裁剪组离中心更远”这一解释。不过同一半径下两组坐标的方向构成不同，而且没有逐样本控制实际可见前景比例，因此仍不能把差值当作裁剪的纯因果效应。
  \item \textbf{位置编码结论只覆盖受控的小型 ViT。}Learnable、2-D Sin-Cos 和 No-Pos 的比较固定为 Center 训练与 CLS 聚合，适合回答当前结构中的位置编码影响，但不能直接外推到 Mean Pooling、卷积 stem 或更大模型。
  \item \textbf{早停依据偏向中心。}验证集使用 Center 协议，可能优先选择中心准确率较好的检查点，而不是网格或复合扰动最稳定的检查点。
  \item \textbf{训练预算有限。}CPU 配置采用 7 像素 patch、64 维 embedding、2 层 Encoder 和最多 12 个 epoch，足够完成课程项目的比较，但没有探索更大 ViT 或更长训练的上限。
  \item \textbf{Attention 解释有限。}rollout 样本按预测差异筛选，不能代表完整测试分布；注意力高响应也不等同于因果特征重要性。
  \item \textbf{效率证据仍有限。}当前已在同一台 CPU、统一线程数与 batch size 下重复测量前向延迟和吞吐量，但尚未测量峰值内存、训练吞吐量与 FLOPs。该基准适合比较本次运行中的相对开销，不用于宣称跨硬件的最佳计算效率。
\end{enumerate}

\subsection{后续可开展实验}

如果继续推进，可以先把两个补充消融从三个随机种子扩展到与主实验一致的五个种子，再在相同增强下单独切换 CLS 与 Mean Pooling、比较有无卷积 stem，形成更完整的因子设计。这样能够在不改变训练预算的前提下，更稳健地估计增强、聚合方式与局部归纳偏置各自的贡献。

对于数据协议，后续可以设计可见像素比例受控的配对样本：不仅保持位移距离相同，还匹配移动方向，并分别保留完整前景和模拟裁剪，从而比当前等半径分组更严格地区分纯位置变化与信息缺失。验证阶段也可同时监控 Center Accuracy 与扰动验证准确率，或者使用二者加权指标选择检查点。完成这些对照后，再考虑 GPU 配置、更小 patch 和更长训练。
